\lecture{7.3}{Double Robustness}
\label{lec:7-3}

\subsection{Large-Sample Limits of AIPTW}

Double robustness is a large-sample property. Suppose
\[
\Qbar_n(a,w)\overset{p}{\longrightarrow}\Qbar^*(a,w),
\qquad
g_{n,a}(w)\overset{p}{\longrightarrow}g_a^*(w),
\]
where $\Qbar^*$ and $g_a^*$ may or may not equal the true nuisance functions
$\Qbar$ and $g_a$. For treatment level $1$, the AIPTW estimator is expected to
converge to
\[
H(\Qbar^*,g_1^*)
=
\E\left[
\Qbar^*(1,W)
+
\frac{\ind(A=1)}{g_1^*(W)}
\{Y-\Qbar^*(1,W)\}
\right].
\]
The question is when this limit equals the target $\E\{\Yone\}$.

\subsection{Correct Outcome Regression}

First suppose the outcome regression is consistently estimated, so that
\[
\Qbar^*(1,W)=\Qbar(1,W)=\E(Y\given A=1,W).
\]
Then
\[
\begin{aligned}
H(\Qbar,g_1^*)
&=
\E\left[
\Qbar(1,W)
+
\frac{\ind(A=1)}{g_1^*(W)}
\{Y-\Qbar(1,W)\}
\right]\\
&=
\E\left[
\Qbar(1,W)
+
\E\left\{
\frac{\ind(A=1)}{g_1^*(W)}
\{Y-\Qbar(1,W)\}
\given W
\right\}
\right].
\end{aligned}
\]
The inner expectation is zero because
\[
\begin{aligned}
\E\left[
\ind(A=1)\{Y-\Qbar(1,W)\}
\given W
\right]
&=
\Prob(A=1\given W)
\E\{Y-\Qbar(1,W)\given A=1,W\}\\
&=
g_1(W)\{ \Qbar(1,W)-\Qbar(1,W)\}\\
&=0.
\end{aligned}
\]
Therefore,
\[
H(\Qbar,g_1^*)=\E\{\Qbar(1,W)\}=\E\{\Yone\}.
\]
The propensity score limit may be wrong, but the estimator is still consistent if the
outcome regression is correct.

\subsection{Correct Propensity Score}

Now suppose the propensity score is consistently estimated, so
$g_1^*(W)=g_1(W)$, but $\Qbar^*$ may be wrong. Then
\[
\begin{aligned}
H(\Qbar^*,g_1)
&=
\E\left[
\Qbar^*(1,W)
+
\frac{\ind(A=1)}{g_1(W)}
\{Y-\Qbar^*(1,W)\}
\right]\\
&=
\E\left[
\Qbar^*(1,W)
+
\E\left\{
\frac{\ind(A=1)}{g_1(W)}
\{Y-\Qbar^*(1,W)\}
\given W
\right\}
\right].
\end{aligned}
\]
Compute the inner expectation:
\[
\begin{aligned}
\E\left[
\frac{\ind(A=1)}{g_1(W)}
\{Y-\Qbar^*(1,W)\}
\given W
\right]
&=
\frac{1}{g_1(W)}
\E\left[
\ind(A=1)\{Y-\Qbar^*(1,W)\}
\given W
\right]\\
&=
\frac{g_1(W)}{g_1(W)}
\{\Qbar(1,W)-\Qbar^*(1,W)\}\\
&=
\Qbar(1,W)-\Qbar^*(1,W).
\end{aligned}
\]
Substituting back,
\[
H(\Qbar^*,g_1)
=
\E\left[
\Qbar^*(1,W)+\Qbar(1,W)-\Qbar^*(1,W)
\right]
=
\E\{\Qbar(1,W)\}
=
\E\{\Yone\}.
\]

\subsection{Two Chances to Be Consistent}

\begin{definition}[Double robustness]
An estimator is doubly robust for a parameter if it is consistent when either of two
nuisance components is consistently estimated, even if the other nuisance component
is inconsistently estimated.
\end{definition}

For AIPTW, the two nuisance components are $\Qbar$ and $g$. The estimator of
$\E\{\Yone\}$ is consistent if either $\Qbar_n(1,w)$ is consistent for $\Qbar(1,w)$ or
$g_{n,1}(w)$ is consistent for $g_1(w)$. The same reasoning applies to
$\E\{\Yzero\}$ and therefore to the ATE.

\begin{identificationbox}
Under the causal identification assumptions, the AIPTW estimator targets
$\E\{\Ya{a}\}$. Its large-sample limit equals this target if either
$\Qbar_n(a,w)$ consistently estimates $\Qbar(a,w)$ or $g_{n,a}(w)$ consistently
estimates $g_a(w)$.
\end{identificationbox}

\subsection{Concrete Misspecification Example}

A simple data-generating system makes the double-robustness statement
operational. Suppose
\[
  W\sim \operatorname{Uniform}(-2,2),
  \qquad
  A\given W=w\sim \operatorname{Bernoulli}\{g_0(w)\},
\]
where
\[
  g_0(w)=\operatorname{expit}\left\{\frac{3}{2}(w+1)^2-3\right\},
\]
and suppose the conditional mean outcome is
\[
  \Qbar_0(a,w)=\E(Y\given A=a,W=w)=1+a-w-aw.
\]
The true outcome regression contains an interaction between treatment and the
covariate because of the term $-aw$. A main-terms-only outcome regression, such
as $\E(Y\given A,W)=\beta_0+\beta_1A+\beta_2W$, is therefore misspecified. The
true propensity score is nonlinear in $W$ because it contains $(w+1)^2$. A
logistic propensity score model with only a linear term in $W$ is therefore
misspecified.

This example creates four useful cases:
\[
\begin{array}{c|c|c|c}
\text{Outcome model} & \text{Propensity model} & \text{G-computation}
& \text{IPTW}\\
\hline
\text{correct} & \text{correct} & \text{consistent} & \text{consistent}\\
\text{correct} & \text{wrong} & \text{consistent} & \text{biased in general}\\
\text{wrong} & \text{correct} & \text{biased in general} & \text{consistent}\\
\text{wrong} & \text{wrong} & \text{biased in general} & \text{biased in general}
\end{array}
\]
The AIPTW estimator is consistent in the first three rows, but not generally in
the fourth row, where both nuisance models are wrong. In the second row, the augmentation term uses
the correct outcome regression and the residual correction has mean zero. In the
third row, the residual correction is weighted by the correct propensity score
and removes the bias left by the wrong outcome regression. The fourth row is the
warning case: double robustness gives two chances to remove bias, not a
guarantee when both chances fail.

\subsection{Why Double Robustness Is Not a License for Bad Models}

Double robustness is sometimes summarized as ``two chances to get it right.'' That
phrase is useful, but it can be misleading. If both nuisance estimators are poor, the
AIPTW estimator can be badly biased. If both are estimated by simplistic parametric
models that miss important nonlinearities or interactions, there may be no real second
chance.

Double robustness also does not protect against violations of causal assumptions.
Unmeasured confounding, interference, ill-defined interventions, or severe
positivity violations can invalidate the causal interpretation regardless of the
statistical estimator.

\subsection{Connection to Efficiency}

The AIPTW estimator is closely related to the efficient influence function for the
counterfactual mean. When both nuisance functions are consistently estimated at
sufficient rates, AIPTW can achieve the nonparametric efficiency bound for the target
parameter. This means that, in large samples and under regularity conditions, no
regular asymptotically linear estimator can have smaller asymptotic variance in the
nonparametric model.

\begin{keyidea}
Double robustness concerns consistency under one correct nuisance function.
Efficiency concerns variance when both nuisance functions are sufficiently well
estimated. The same augmentation structure is responsible for both properties.
\end{keyidea}
